\section{Mathematical Details}
\label{app:math-details}

\paragraph{Route features and observation.}  Each candidate route is
represented by mean pooling its segment latents, and the shared context
summarizes the graph and decision progress:
\begin{equation}
\begin{array}{rcl}
\varphi_{i,t,k}
&=&
\displaystyle
\frac{1}{|P^k_{i,t}|}
\sum_{u\in P^k_{i,t}}z_{t,u}.
\end{array}
\end{equation}
\begin{equation}
\begin{array}{rcl}
\bar z_t
&=&
\displaystyle
\frac{1}{|N|}\sum_{u\in N}z_{t,u},\\[0.35em]
g_t
&=&
\displaystyle
\left[
\bar z_t,
\frac{t}{T},
\frac{|\mathcal{I}_t|}{|N^{\mathrm{CAV}}|}
\right].
\end{array}
\end{equation}
The vectorized realization of the observation operator is therefore
\begin{equation}
o_{i,t}
=
\big[
\varphi_{i,t,1}\Vert\cdots\Vert
\varphi_{i,t,K}\Vert g_t
\big],
\end{equation}
where \(|N^{\mathrm{CAV}}|\) is the number of CAVs scheduled in the episode
and \(\Vert\) denotes concatenation.

\paragraph{Advantage normalization and PPO surrogate.}  The combined
advantage is standardized over minibatch statistics \(\mu_A,\sigma_A\) and
clipped to \([-c,c]\), and the probability ratio is clipped to
\([1-\epsilon,1+\epsilon]\):
\begin{equation}
\begin{array}{rcl}
\bar A_{i,t}^{\mathrm{actor}}
&=&
\displaystyle
\mathrm{clip}\!\left(
\frac{A_{i,t}^{\mathrm{actor}}-\mu_A}{\sigma_A+\varepsilon_A},-c,c
\right),\\[0.35em]
\rho_{i,t}
&=&
\displaystyle
\frac{\pi_\theta(a_{i,t}\mid o_{i,t})}
{\pi_{\mathrm{old}}(a_{i,t}\mid o_{i,t})},\\[0.35em]
\bar\rho_{i,t}
&=&
\mathrm{clip}(\rho_{i,t},1-\epsilon,1+\epsilon).
\end{array}
\end{equation}
The actor minimizes the clipped surrogate with entropy regularization:
\begin{equation}
\begin{array}{rcl}
\mathcal{L}_{\pi}
&=&
-\mathbb{E}_{i,t}
\left[
\min\!\left(
\rho_{i,t}\bar A_{i,t}^{\mathrm{actor}},
\bar\rho_{i,t}\bar A_{i,t}^{\mathrm{actor}}
\right)
\right]\\
&&
-\beta\mathbb{E}_{i,t}
\left[\mathcal{H}(\pi_\theta(\cdot\mid o_{i,t}))\right],
\end{array}
\end{equation}
Here \(\mathbb{E}_{i,t}\) averages over active CAV decisions in the imagined
minibatch, \(\mathcal H\) is the entropy of the categorical route
distribution, and \(\beta\) is its weight.

\paragraph{Per-vehicle lambda return.}  Returns are linked by vehicle
identity.  Let \(\zeta_{i,t+1}^{\mathrm{act}}\in\{0,1\}\) indicate
whether CAV \(i\) remains in \(\mathcal{I}_{t+1}^{\mathrm{im}}\).  The
masked bootstrap value, per-vehicle lambda return, and critic loss are
\begin{equation}
\begin{array}{rcl}
\widetilde v_{i,t+1}^{\mathrm{veh}}
&=&
\zeta_{i,t+1}^{\mathrm{act}}v_{i,t+1}^{\mathrm{veh}},\\
G_{i,t}^{\lambda}
&=&
r_{i,t}^{\mathrm{veh}}
+\gamma\left[
(1-\lambda)\widetilde v_{i,t+1}^{\mathrm{veh}}
+\lambda G_{i,t+1}^{\lambda}
\right],\\[0.35em]
\mathcal{L}_{V}^{\mathrm{veh}}
&=&
\displaystyle
\frac{1}{2}\mathbb{E}_{i,t}
\left[
(v_{i,t}^{\mathrm{veh}}
-\mathrm{sg}(G_{i,t}^{\lambda}))^2
\right].
\end{array}
\end{equation}
Only observations of the same CAV are connected by this recursion.  Once the
vehicle leaves the active set, \(\zeta_{i,t+1}^{\mathrm{act}}=0\) prevents
bootstrapping beyond its imagined lifetime.

\paragraph{Global lambda return.}  The global critic is trained from the
corresponding epoch-level lambda return:
\begin{equation}
\begin{array}{rcl}
G_t^{\mathrm{global}}
&=&
r_t^{\mathrm{global}}
+\gamma\left[
(1-\lambda)v_{t+1}^{\mathrm{global}}
+\lambda G_{t+1}^{\mathrm{global}}
\right],\\[0.35em]
\mathcal{L}_{V}^{\mathrm{global}}
&=&
\displaystyle
\frac{1}{2}\mathbb{E}_{t}
\left[
(v_t^{\mathrm{global}}
-\mathrm{sg}(G_t^{\mathrm{global}}))^2
\right].
\end{array}
\end{equation}
Here \(\mathbb{E}_t\) averages over imagined decision epochs.


\paragraph{GNN traffic encoder.}  Given the traffic field \(X_t\) and the
row-normalized adjacency matrix \(A\), the encoder applies
\begin{equation}
\begin{array}{rcl}
H^0 &=& X_tW_x,\\
H^{\ell+1} &=&
\mathrm{LN}\!\left(
H^\ell+\sigma(AH^\ell W_m+H^\ell W_s)
\right)
\end{array}
\end{equation}
where \(W_x,W_m,W_s\) are learned projections, \(\sigma\) is a nonlinear
activation, and \(\mathrm{LN}\) denotes layer normalization.

\paragraph{Decoder reconstruction loss.}  The decoder is trained jointly with
the dynamics model with a reconstruction loss on observed traffic fields:
\begin{equation}
\begin{array}{rcl}
\mathcal{L}_{\mathrm{rec}}
&=&
\displaystyle
\frac{1}{M}\sum_{u\in N}
\|\hat x^{\mathrm{dyn}}_{t,u}-x^{\mathrm{dyn}}_{t,u}\|^2.
\end{array}
\end{equation}

\paragraph{Weighted flow-matching loss.}  Because most road segments are
empty most of the time, the training loss applies per-segment weights
\(w_{t,u}\) that emphasize occupied and congested segments, with a bounded
maximum weight to prevent degenerate focusing:
\begin{equation}
\begin{array}{l}
\mathcal{L}_{\mathrm{FM}}
=
{E}_{t,\tau}
\left[
\sum_{u\in N}w_{t,u}\|\hat\nu_{t,\tau,u}-\nu_{t,\tau,u}\|^2
\right],
\end{array}
\end{equation}

\paragraph{HDV response target and prediction.}  For an observed episode of
length \(T\), the supervision target aggregates the observed HDV background
field \(h^{\mathrm{HDV}}_t\) over time:
\begin{equation}
\begin{array}{rcl}
H^\star
&=&
\mathrm{Norm}\!\left(
\displaystyle\frac{1}{T}\sum_{t=1}^{T} h^{\mathrm{HDV}}_t
\right),
\end{array}
\end{equation}
where \(\mathrm{Norm}\) clips negative entries and rescales the vector to
sum to one.  The predictor summarizes the previous episode's CAV pressure and
estimates the next HDV response field with an MLP followed by a softmax:
\begin{equation}
\begin{array}{rcl}
\bar\Psi_{e-1}^{\mathrm{CAV}}
&=&
\mathrm{Norm}\!\left(
\displaystyle\frac{1}{T_{e-1}}
\sum_{t=0}^{T_{e-1}-1}\Psi_{e-1,t}^{\mathrm{CAV}}
\right),
\end{array}
\end{equation}
\begin{equation}
\begin{array}{rcl}
\hat H_e^{\mathrm{HDV}}
&=&
g_\psi^{\mathrm{HDV}}
\!\left(
\bar\Psi_{e-1}^{\mathrm{CAV}},
H_{e-1}^{\mathrm{HDV}}
\right),\\[0.35em]
g_\psi^{\mathrm{HDV}}(\cdot)
&=&
\mathrm{softmax}\!\left(f_\psi^{\mathrm{HDV}}(\cdot)\right).
\end{array}
\end{equation}
The response loss is a soft-label cross entropy over adjacent episode pairs,
and the response predictor and the base world model use separate optimizers:
\begin{equation}
\begin{array}{rcl}
\mathcal{L}_{\mathrm{resp}}
&=&
\displaystyle
-
\frac{1}{|\mathcal{E}_{\mathrm{pair}}|}
\sum_{e\in\mathcal{E}_{\mathrm{pair}}}
\sum_{u\in N}
H_e^\star(u)
\log \hat H_e^{\mathrm{HDV}}(u).
\end{array}
\end{equation}
\begin{equation}
\begin{array}{rcl}
\mathcal{L}_{\phi}
&=&
\mathcal{L}_{\mathrm{FM}}
+
\lambda_{\mathrm{rec}}\mathcal{L}_{\mathrm{rec}},\\[0.35em]
\mathcal{L}_{\psi}
&=&
\mathcal{L}_{\mathrm{resp}}.
\end{array}
\end{equation}
Here \(\mathcal{E}_{\mathrm{pair}}\) is the set of training episodes whose
previous episode is available.  \(\mathcal{L}_{\phi}\) updates the encoder,
flow dynamics, and decoder, while \(\mathcal{L}_{\psi}\) updates only the
response predictor \(g_\psi^{\mathrm{HDV}}\).

\paragraph{Route candidate generation.}  Candidate routes are generated by
repeated Dijkstra search on the segment graph.  The initial segment cost is
the free-flow travel time \(c^1(u)=\ell(u)/v_{\max}(u)\).  After a route is
found, all segment nodes on it are penalized before the next search:
\begin{equation}
\begin{array}{rcl}
P^k_{i,t}
&=&
\mathrm{Dijkstra}(G,u_{i,t},g_i;c^k),\\
c^{k+1}(u)
&=&
\left\{
\begin{array}{ll}
\eta c^k(u), & u\in P^k_{i,t},\\
c^k(u), & u\notin P^k_{i,t},
\end{array}
\right.
\end{array}
\end{equation}
where \(\eta>1\) is the route-diversity penalty.  The search stops when no
new valid route can be found.
